\section{Related Work}
\label{sec:related}

\subsection{Open-world and zero-shot NER}
The three open-world systems we compare against differ in backbone as much as in approach. GLiNER \citep{zaratiana2024gliner} scores spans against candidate type names with a DeBERTaV3 encoder \citep{he2023debertav3}. GNER \citep{ding2024gner} generates them with a sequence-to-sequence model \citep{raffel2020t5}. OWNER \citep{genest2025owner} trains a dedicated entity encoder, then clusters mentions into dynamically named types. Typing from type descriptions \citep{aly2021zeroshot} needs no labelled example at all. A complementary line distils or prompts billion-parameter language models for open NER, as in UniversalNER \citep{zhou2024universalner}, GoLLIE \citep{sainz2024gollie} and ChatIE \citep{wei2024chatie}. Our comparison also probes that compute and energy cost, \rev{through an instruction-tuned LLM baseline (Qwen2.5-1.5B in $4$-bit) sized to the hardware this paper measures everything on}. OPENER differs from all of these by reusing a pretrained embedding rather than training an encoder, which is what keeps it cheap enough for the legal setting we target.

\subsection{Legal NLP and legal NER}
Legal NLP is an active field, spanning summarisation, judgement prediction, question answering and information extraction, as a recent legal-NLP survey documents \citep{ariai2025legalnlp}. It is driven by the specialised vocabulary, the long formulaic documents and the jurisdiction-specific entity inventories of legal text. On the modelling side, domain-adapted language models such as LEGAL-BERT \citep{chalkidis2020legalbert} and multi-task benchmarks such as LexGLUE \citep{chalkidis2022lexglue} show that in-domain pretraining and evaluation matter for legal NLU.

The dominant approach to legal NER is to train or adapt a supervised tagger for each corpus, schema and language. LegNER \citep{karamitsos2025legner} domain-adapts BERT with span supervision for English legal NER and anonymisation. A Serbian legal NER study \citep{kalusev2025serbian} fine-tunes a BERT-type model for court rulings and releases a new dataset. The LegalEval shared task (SemEval-2023 Task~6) spurred systems such as a Bi-LSTM with Flair embeddings for Indian court judgements \citep{ramesh2023flairnlp}. Even LLM-based studies remain schema-bound. A financial NER study \citep{lu2025financial} prompts GPT-4o, LLaMA-3.1 and Gemini-1.5, and reports systematic failure modes on domain entities. These systems achieve strong in-corpus scores, but each one requires labelled data and (re)training tied to one schema and one language. E-NER \citep{au2022ener} annotates U.S.\ SEC/EDGAR filings with seven classes, including court, government body and legislation/act. \rev{Indian-Legal \citep{kalamkar2022indianlegal} tags Indian judgements with $14$ types, from judge and lawyer to precedent.} LeNER-Br \citep{luzdearaujo2018lenerbr} covers Brazilian court decisions with dedicated \emph{legislation} and \emph{case-law} tags. German-LER \citep{leitner2020germanler} annotates German federal decisions with a fine-grained 19-class schema.

Crucially, the E-NER study \citep{au2022ener} reports that NER models trained on general English degrade sharply on legal text, \rev{from $90.5$ F1 in-domain down to $61.1$ once tested on legal filings}. That is precisely the distribution shift we study. Rather than retraining a legal-specific tagger per corpus and language, we build a single open-world pipeline that no legal schema ever touches, and we measure how far it goes on all four of them (\autoref{sec:experiments}).

\subsection{Energy-aware evaluation}
Reporting efficiency alongside accuracy is increasingly advocated, from the Green AI agenda \citep{schwartz2020greenai} to energy audits of NLP training \citep{strubell2019energy}, systematic carbon reporting \citep{henderson2020systematic,patterson2022carbon} and whole-lifecycle views of AI's footprint \citep{wu2022sustainable}. At inference time specifically, a deployment-cost study \citep{luccioni2024power} finds general-purpose generative models orders of magnitude more expensive per query than task-specific ones. That is precisely the gap this frugal pipeline targets. We adopt the same three-axis view, on quality, latency and energy, measured with CodeCarbon \citep{codecarbon2024} and reported in the energy-times-grid-intensity accounting of \citet{lacoste2019quantifying}.
